\chapter*{Abstract}
\addcontentsline{toc}{chapter}{Abstract}
Left truncation occurs when individuals enter observation after the origin of
the survival time scale and are observed only if they survive to entry. This
thesis extends the \pkg{} R package to left-truncated, right-censored,
single-event survival data. Subject-specific entry times are passed through
the R and C interfaces and used in default log-rank splitting, terminal-node
estimation, and selected performance calculations. Regression checks address
ordinary-data compatibility, saved-forest restoration, row alignment, tuning,
and delayed-entry scoring.

The completed simulation study uses six data-generating processes and
\SimReplications{} paired replications per scenario. Each replication contains
\SimTrainN{} retained training subjects, \SimTrees{} trees per forest, and
\SimTestN{} independently drawn population test covariates. Entry-aware fitting
reduces mean absolute survival-curve error in every scenario: by
16.66--91.39\% for Cox regression and 31.05--66.01\% for random survival forests.
Entered-cohort Brier/CRPS and time-averaged AUC error also decrease in all six
scenarios. Concordance changes are much smaller and include deterioration for
RSF in some scenarios. These metrics address different prediction targets and
use different evaluation samples across model families, so the findings do
not establish universal superiority of one method or one scoring procedure.

The contribution is an applied package extension with reproducible numerical
evaluation. Remaining work includes review of the observation-conditioned
Brier interpretation, the added AUC definition, censoring-model assumptions,
and the scope of a future upstream contribution.

\draftnote{Updated discussion draft using the completed 100-replication study.
The literature review, additional evaluation requested in the project outline,
and institutional formatting remain unfinished. Numerical results use package
commit \RunCommit{}; no upstream submission is claimed.}
